\section{Introduction}
\label{sec:intro}

Distance to the nearest hospital is one of the most widely used
measures of healthcare access in applied economics. It motivates the
literature on hospital closures \citep{avdic2016, gujralBasu2019, carroll2019},
on rural healthcare deserts \citep{germack2019, lindrooth2018}, and
on the spatial allocation of health resources \citep{gaynorTown2011}. Empirical work uses haversine kilometers, road
travel time \citep{doyle2011}, or driving time over OpenStreetMap as
inputs to demand systems, instrumental variables for utilization, or
exposure measures in event studies of facility entry and exit. Yet
geographic distance is, at best, a noisy proxy for the medically
meaningful concept of \emph{network access}---the actual ease with which
patients reach the appropriate level of care.

This gap between geography and access is particularly stark in Brazil,
where the public Sistema Único de Saúde (SUS) covers $\approx 75\%$ of
the population through a hub-and-spoke regionalization built atop a
$5{,}570$-municipality administrative grid. Many patients in Amazonian
municipalities reach high-complexity care exclusively by river boat or
chartered plane; many in the semi-arid Northeast bypass nominally
nearest hospitals because those facilities lack qualified specialists,
ICU beds, or suitable equipment. In both cases, the metric ``kilometers
to the nearest hospital'' fundamentally misses the structure of
realized access. The error is not random: it is concentrated in the
poorest and most epidemiologically vulnerable regions of the country.

\paragraph{This paper.}
We provide a graph-based, data-driven measure of healthcare access for
the universe of Brazilian municipalities, and use it to map where the
geographic and the network-revealed views of access agree and where
they sharply diverge. Our object of analysis is the bipartite patient
flow graph linking $5{,}588$ municipalities of residence to $5{,}828$
public hospitals through $\approx 100$ million federal hospital
admissions (AIH) over $2010\text{--}2024$, the entire universe of SUS
inpatient activity in the period. We learn a $128$-dimensional
representation of each node using a sparse implementation of node2vec
\citep{grover2016} that encodes second-order random walks with
edge-weight bias. Two embeddings are produced: a bipartite embedding
in which municipalities and hospitals share a common space (allowing
direct municipality-to-hospital cosine distance), and a TF--IDF-weighted
co-patiency projection of the bipartite graph onto the
municipality-only space.

For each municipality, we compute two scalar measures of isolation
relative to the set of high-complexity hospitals---those with active
SUS habilitation in cardiology, oncology, neurology, nephrology,
perinatology, trauma-orthopedics, or transplantation:
\[
\isokm_i = \text{haversine km to nearest high-complexity hospital},\qquad
\isoemb_i = 1 - \cos(\mathbf{e}_i, \mathbf{e}_{h^*(i)}),
\]
Standardizing each to a $z$-score and taking the difference yields the
\emph{divergence}
\(
\divz_i = z(\isoemb_i) - z(\isokm_i).
\)
A positive divergence indicates that the embedding sees a municipality
as \emph{more} isolated than the map does---patients leapfrog the
nominally nearest hospital because it is inadequate. A negative
divergence indicates the opposite: the geography overstates isolation
because patients reach a distant hub via channels (river, air, informal
referral pathways across state lines) that road geography ignores.

\paragraph{Findings.}
Three patterns emerge.

First, the divergence is large. The cross-sectional standard deviation
of $\divz$ is $0.91$, with $9\%$ of municipalities exhibiting
$|\divz| > 1.5$. Roughly $56.7\%$ of municipalities have a different
nearest high-complexity hospital under the embedding metric than under
the geographic metric.

Second, the divergence is geographically organized into recognizable
macro-regional patterns. Northern (Amazonian) Brazil is dominantly
\emph{negatively} divergent: $\bar{\divz}_{\text{N}} = -0.22$, with
median geographic distance to the nearest high-complexity hospital of
$30$\,km but embedding-revealed connections that pull these municipalities
toward Manaus, Boa Vista, and Belém despite distances often exceeding
$200$\,km by air or river. Northeastern semi-arid municipalities show
the opposite signature in pockets of Piauí, Maranhão, and Bahia interior:
nominal hospital proximity (medians $\approx 19$\,km) coexists with high
embedding-isolation, as patients route around inadequate local
facilities to reach state-capital tertiary care. The Southeast and
Center-West show positively-skewed but moderate divergence, consistent
with metropolitan saturation effects in Greater São Paulo and the
Brasília corridor. The South emerges as the most well-calibrated region:
geographic and network-revealed access are closely aligned, with
$\bar{\divz}_{\text{S}} = -0.10$ and the lowest cross-sectional dispersion.

Third, despite the strong descriptive signal, we cannot detect a
robust cross-sectional association between divergence and amenable
mortality once standard controls (state fixed effects, log population,
log GDP) are included. Adding our network isolation measure to a
ridge regression already containing $\isokm$ produces a $\Delta R^2$
of $-0.0002$, statistically indistinguishable from zero. We argue this
\emph{null} reflects two well-documented confounders of vital
registration in Brazil: (i) under-reporting of deaths due to
ill-defined causes (R-codes) in the Northern and rural Northeastern
SIM strata exactly where divergence is largest, and (ii) cross-border
death recording, in which patients referred from semi-arid municipalities
die at the receiving capital and have their death registered in the
hospital's municipality rather than the residential one. Both biases
attenuate the mechanical link between access deficit and observed
mortality at the residential-municipality level. We treat the null as
a methodological warning---one that, to our knowledge, the existing
hospital-access literature using SIM aggregates has not articulated---
rather than as evidence that network access does not matter for
mortality, and outline in the discussion a closure event-study design
that sidesteps the confounders by relying on within-municipality
variation.

\paragraph{Contributions and related literature.}
We contribute to three strands of literature.

First, on \emph{hospital closures and outcomes}, the existing literature
\citep{avdic2016, gujralBasu2019, carroll2019, joynt2015} estimates
the average effect of facility loss on local mortality, typically using
distance-to-nearest as the treatment intensity. Our embedding offers a
more granular treatment metric---revealed exposure---that we will exploit
in a companion event-study (Pivot B in our research roadmap, beyond the
scope of the current draft).

Second, on \emph{healthcare deserts and access measurement}
\citep{germack2019, gaynorVogt2003}, we provide the
first nationwide application of patient-flow embeddings to define
hospital deserts, and propose the divergence metric as a falsification
device for any access measure based purely on physical distance. The
$56.7\%$ disagreement rate between geographic and network-revealed
nearest hospital is, to our knowledge, the largest such estimate documented.

Third, on \emph{representation learning for economic networks}
\citep{kelly2021, acemoglu2024, bonhomme2019}, we show that node2vec
applied to a tractable, mid-sized bipartite graph ($\approx 11{,}000$
nodes, $\approx 500{,}000$ edges) yields embeddings that pass strong
sanity checks (top-10 nearest neighbors of São Paulo recover the entire
metropolitan ABC; top-10 of Manaus recover its peri-urban ribeirinho
catchment) while retaining a useful $0.30$ Pearson correlation with
geographic distance---enough to track real spatial structure without
collapsing into a redundant geography.

\paragraph{Roadmap.}
Section \ref{sec:setting} describes the Brazilian institutional
setting. Section \ref{sec:data} introduces the data: SIH-AIH for hospital
flows, SIM for mortality, CNES habilitations for hospital quality, and
IBGE for spatial overlays. Section \ref{sec:method} formalizes the
graph construction, the node2vec specification, and the divergence
metric. Section \ref{sec:results} presents the headline divergence map,
the regional decomposition, the canonical-case taxonomy, and the
mortality null with its diagnosis. Section \ref{sec:discussion}
discusses policy implications and the planned causal extension.
Section \ref{sec:conclusion} concludes.
